\chapter*{Реферат}


Отчет 64 с., 17 рис., 7 табл., 20 источников.

БЕСПЛАТФОРМЕННАЯ ИНЕРЦИАЛЬНАЯ НАВИГАЦИОННАЯ СИСТЕМА, БИНС, МИКРОМЕХАНИЧЕСКИЕ ДАТЧИКИ, МЭМС, ПУТЕИЗМЕРИТЕЛЬНАЯ ТЕЛЕЖКА, КОНТРОЛЬ ГЕОМЕТРИИ ЖЕЛЕЗНОДОРОЖНОГО ПУТИ, STM32, КВАТЕРНИОНЫ, ФИЛЬТРАЦИЯ СИГНАЛОВ, ТЕМПЕРАТУРНАЯ КОМПЕНСАЦИЯ.

Объект исследования --- процессы контроля геометрических параметров железнодорожного пути путеизмерительными тележками.

Предмет исследования --- алгоритмы и аппаратные решения бесплатформенной инерциальной навигационной системы на микромеханических датчиках для путеизмерительной тележки.

Цель работы --- разработка бесплатформенной инерциальной навигационной системы на основе низкоточных микромеханических датчиков для путеизмерительной тележки, обеспечивающей контроль геометрических параметров железнодорожного пути при снижении себестоимости и снаряжённой массы по сравнению с существующими решениями.

Методы исследования: системный анализ путеизмерительных средств; статистическая обработка сигналов методом Уэлфорда с отбраковкой выбросов по критерию $2\sigma$ и кусочно-линейная температурная компенсация смещения нуля; описание ориентации в аппарате кватернионов; численное интегрирование показаний датчиков; натурные испытания макета.

В результате работы реализованы прошивка микроконтроллера STM32F3 для сбора и первичной обработки данных микромеханических датчиков, навигационный алгоритм счисления координат и ориентации, программное обеспечение персонального компьютера, печатная плата коммутации и защитный корпус прибора. Себестоимость инерциального измерительного модуля снижена с нескольких миллионов до десятков тысяч рублей; результаты работы использованы при разработке перспективной модели путеизмерительной тележки в ЗАО \enquote{ПИК Прогресс}.

\newpage